% !TeX root = ../抽象代数笔记.tex
% 由主文档通过 \input 载入；请编译主文档。
\setcounter{chapter}{18} % 保留原书章号。
\chapter{根式扩张与方程的可解性}\label{ch:radicals}

\begin{description}
  \item[学习目标] 理解根式表达与可解群之间的联系，准确陈述一般高次方程不可用根式求解的含义。
  \item[先修知识] 第\ref{ch:solvable-groups}、\ref{ch:galois}章。本章基域均假设特征为 $0$。
\end{description}

\section{根式扩张与根式可解的定义}
定义根式扩张（radical extension）为可由有限域塔得到的扩张，其中每一步形如 $F_i=F_{i-1}(\alpha_i)$，并且存在正整数 $m_i$ 使 $\alpha_i^{m_i}\in F_{i-1}$。定义多项式根式可解，是指其全部根包含于某个根式扩张中；说明分裂域本身不必就是所选的根式扩张。

\section{单位根与圆分扩张}
定义本原 $n$ 次单位根 $\zeta_n$ 和圆分多项式（cyclotomic polynomial）$\Phi_n(x)$。证明 $\Phi_n\in\mathbb Z[x]$ 且在 $\mathbb Q[x]$ 中不可约，从而建立 $\operatorname{Gal}(\mathbb Q(\zeta_n)/\mathbb Q)\cong(\mathbb Z/n\mathbb Z)^\times$。解释添加单位根为何能控制纯根式方程的自同构。

\section{根式可解与可解群的等价}
设 $F$ 为特征 $0$ 的域，$f\in F[x]$ 为非恒定多项式，$L$ 为其分裂域。证明 $f$ 在 $F$ 上根式可解当且仅当 $\operatorname{Gal}(L/F)$ 为可解群。分别安排两个方向：对根式塔添加必要单位根并控制其正规闭包；对可解群逐层取固定域，并利用含单位根时的素数阶循环扩张构造根式。后一个步骤需单独证明，不能仅从“商群阿贝尔”直接跳到“能开根”。

\section{低次方程与一般五次方程}
结合 $S_2,S_3,S_4$ 的可解性解释二、三、四次方程存在根式解；用 $S_n$（$n\geq5$）的不可解性分析一般高次方程。若表述“一般 $n$ 次方程”，明确以代数独立系数上的有理函数域为基域，并证明其伽罗瓦群为 $S_n$。同时安排一个伽罗瓦群为 $S_5$ 的具体有理系数五次多项式，完整验证不可约性与群的判定。

\section{例题、练习与自检}
比较 $x^5-2$ 与不可根式求解的五次多项式，说明次数相同不决定根式可解性。回访尺规作图：以根式塔中的二次步骤解释作图限制，将正多边形判别列为选学。练习要求区分无通用根式公式、某个方程无根式解、方程没有根及无法数值近似这四个不同命题。

